\documentclass[8pt]{article}
\usepackage[utf8]{inputenc}
\usepackage[T2A]{fontenc}
\usepackage[russian]{babel}
\usepackage{amsmath,amssymb}
\usepackage{geometry}
\geometry{a4paper, margin=0.5cm}
\pagestyle{empty}

% Сжатие отступов для размещения на одной странице
\setlength{\parindent}{0pt}
\setlength{\parskip}{0pt plus 1pt}
\setlength{\abovedisplayskip}{2pt}
\setlength{\belowdisplayskip}{2pt}
\setlength{\abovedisplayshortskip}{2pt}
\setlength{\belowdisplayshortskip}{2pt}
\setlength{\jot}{1pt}
\renewcommand{\baselinestretch}{0.92}
\setlength{\fboxsep}{1.5mm}
\setlength{\fboxrule}{0.4pt}

% Команда для создания прямоугольного блока
\newcommand{\block}[1]{%
  \fbox{\begin{minipage}[t]{\dimexpr0.48\textwidth-2\fboxsep-2\fboxrule\relax}
    \raggedright #1
  \end{minipage}}%
}

\begin{document}
\small

\block{
\subsubsection*{1. Метод решения линейного неоднородного уравнения первого порядка}
Уравнение вида: $y' + a(x)y = b(x)$
\textbf{Алгоритм решения:} Метод вариации постоянной: 1). Решаем соответствующее однородное уравнение $y' + a(x)y = 0$: $\frac{dy}{y} = -a(x)dx \Rightarrow y = Ce^{-\int a(x)dx}$.
2). Заменяем постоянную $C$ на функцию $C(x)$: $y = C(x)e^{-\int a(x)dx}$.
3). Подставляем в исходное уравнение и находим $C'(x)$: $C'(x) = b(x)e^{\int a(x)dx}$.
4). Интегрируем: $C(x) = \int b(x)e^{\int a(x)dx}dx + C$.
5). Общее решение: $y = e^{-\int a(x)dx}\left(\int b(x)e^{\int a(x)dx}dx + C\right)$.
} \hfill \block{
\subsubsection*{2. Уравнение Бернулли}
$y' + a(x)y = b(x)y^m, \quad m \neq 1$
Решение: 1) Делим на $y^m$ (при $y \neq 0$): $y^{-m}y' + a(x)y^{1-m} = b(x)$.
2). Замена $z = y^{1-m}$, тогда $z' = (1-m)y^{-m}y'$.
3). Получаем ЛУ относительно $z$: $\frac{z'}{1-m} + a(x)z = b(x) \qquad z' + (1-m)a(x)z = (1-m)b(x)$.
4). Решаем полученное ЛУ методом вариации постоянной.
5). Возвращаемся к $y$: $y = z^{\frac{1}{1-m}}$.
}

\vspace{2mm}

\block{
\subsubsection*{3. Задача Коши для уравнения $y' = f(x, y)$}
\textbf{Постановка задачи Коши:}
Найти решение $y = \varphi(x)$ дифференциального уравнения $y' = f(x, y)$, удовлетворяющее начальному условию:
$$y(x_0) = y_0$$
\textbf{Теорема существования и единственности:}
Пусть в некоторой области $D$ плоскости $OXY$ содержатся точка $(x_0, y_0)$ и прямоугольник $R = \{|x-x_0| \leq a, |y-y_0| \leq b\}$.
\textbf{а) Условие существования решения:}
Если функция $f(x, y)$ непрерывна в области $D$, то существует решение $y = \varphi(x)$ задачи Коши, определенное в некотором интервале $|x-x_0| < h$.
\textbf{б) Условие единственности решения:}
Если функция $f(x, y)$ непрерывна в $D$ и удовлетворяет условию Липшица по $y$: $|f(x, y_1) - f(x, y_2)| \leq L|y_1 - y_2|$,
где $L > 0$, то решение задачи Коши единственно.
\textbf{Достаточное условие выполнения условия Липшица:}
Если $\frac{\partial f}{\partial y}$ существует и непрерывна в $D$, то $f(x, y)$ удовлетворяет условию Липшица.
} \hfill \block{
\subsubsection*{4. Теорема существования и единственности для системы ОДУ}
\textbf{Постановка задачи:}
Найти решение $\mathbf{y}(x) = (y_1(x), y_2(x), \ldots, y_n(x))^T$ системы $\mathbf{y}' = \mathbf{f}(x, \mathbf{y})$,
удовлетворяющее условию $\mathbf{y}(x_0) = \mathbf{y}_0$.
\textbf{Теорема:}
Пусть в области $D \subset \mathbb{R}^{n+1}$ содержится точка $(x_0, \mathbf{y}_0)$.
Если: 1). Функции $f_i(x, \mathbf{y})$ непрерывны в $D$.
2). Функции $f_i(x, \mathbf{y})$ удовлетворяют условию Липшица по $\mathbf{y}$ в $D$:
$\|\mathbf{f}(x, \mathbf{y}_1) - \mathbf{f}(x, \mathbf{y}_2)\| \leq L\|\mathbf{y}_1 - \mathbf{y}_2\|$, где $L > 0$.
Тогда существует единственное решение $\mathbf{y}(x)$ задачи Коши, определенное в некотором интервале $|x-x_0| < h$.
}

\vspace{2mm}

\block{
\subsubsection*{5. Задача Коши для неявного уравнения}
\textbf{Постановка задачи Коши для уравнения $F(x, y, y') = 0$:}
Найти функцию $y = \varphi(x)$, удовлетворяющую условиям: 1). $F(x, \varphi(x), \varphi'(x)) = 0$ для всех $x$ из некоторого интервала. 2). $\varphi(x_0) = y_0$.
3). $\varphi'(x_0) = y_0'$ (если требуется)
где $y_0'$ определяется из уравнения $F(x_0, y_0, y_0') = 0$.
\textbf{Теорема существования и единственности:}
Если в окрестности точки $(x_0, y_0, y_0')$:
1). $F(x, y, p)$ непрерывна и имеет непрерывные частные производные.
2). $F(x_0, y_0, y_0') = 0$.
3). $\frac{\partial F}{\partial p}(x_0, y_0, y_0') \neq 0$.
Тогда существует единственное решение задачи Коши в некоторой окрестности точки $x_0$.
} \hfill \block{
\subsubsection*{6. Особое решение дифференциального уравнения}
Особым решением дифференциального уравнения называется решение, которое не может быть получено из общего решения ни при каком значении произвольной постоянной.
Особое решение представляет собой огибающую семейства интегральных кривых, определяемых общим решением.
\textbf{Необходимые и достаточные условия:}
Для уравнения $F(x, y, y') = 0$ решение $y = \varphi(x)$ является особым тогда и только тогда, когда:
$$1). \begin{cases}
F(x, y, p) = 0 \\
\frac{\partial F}{\partial p}(x, y, p) = 0, \text{  где } $p = y'$.
\end{cases}$$
$$2). \begin{cases}
\Phi(x, y, C) = 0 \\
\frac{\partial \Phi}{\partial C}(x, y, C) = 0, \text{  где  $\Phi(x, y, C) = 0$ — общее решение.}
\end{cases}$$
Признаки особого решения: 1). В точках особого решения нарушается единственность решения.
2). Через каждую точку особого решения проходит как минимум две интегральные кривые.
3). Особое решение касается каждой интегральной кривой семейства.
}

\vspace{2mm}

\block{
\subsubsection*{7. Простейшая задача вариационного исчисления}
\textbf{Постановка задачи:}
Найти функцию $y(x)$, удовлетворяющую экстремуму функционала: $J[y] = \int_a^b F(x, y, y')dx$ при граничных условиях: $y(a) = y_a, \quad y(b) = y_b$
\textbf{Уравнение Эйлера-Лагранжа:}
Необходимое условие экстремума: $\frac{\partial F}{\partial y} - \frac{d}{dx}\left(\frac{\partial F}{\partial y'}\right) = 0$.
Частные случаи: 1). Если $F = F(x, y')$ (не зависит от $y$): $\frac{\partial F}{\partial y'} = C$. 2). Если $F = F(y, y')$ (не зависит от $x$): $F - y'\frac{\partial F}{\partial y'} = C$
3). Если $F = F(x, y)$ (не зависит от $y'$): $\frac{\partial F}{\partial y} = 0$
Алгоритм решения: 1). Составить уравнение Эйлера-Лагранжа. 2). Решить полученное дифференциальное уравнение. 3). Найти постоянные интегрирования из граничных условий
} \hfill \block{
\subsubsection*{8. Теорема Штурма}
\textbf{Теорема:}
Пусть даны два уравнения: $(p_1(x)y')' + q_1(x)y = 0$ и $(p_2(x)y')' + q_2(x)y = 0$, где $0 < p_2(x) \leq p_1(x)$ и $q_1(x) \leq q_2(x)$ на отрезке $[a, b]$.
Если решение $y_1(x)$ первого уравнения имеет нули в точках $x = a$ и $x = b$, то любое нетривиальное решение $y_2(x)$ второго уравнения имеет хотя бы один нуль на интервале $(a, b)$.
\textbf{Сл. 1 (о разделении нулей):}
Нули двух линейно независимых решений уравнения $y'' + q(x)y = 0$ чередуются, т.е. между любыми двумя соседними нулями одного решения лежит ровно один нуль другого решения.
\textbf{Сл. 2 (о числе нулей):}
Если $q_1(x) < q_2(x)$ на $[a, b]$, то между любыми двумя соседними нулями нетривиального решения уравнения $y'' + q_1(x)y = 0$ лежит хотя бы один нуль любого нетривиального решения уравнения $y'' + q_2(x)y = 0$.
}

\vspace{2mm}

\block{
\subsubsection*{9. Формула Лиувилля-Остроградского для системы ОДУ}
Для линейной однородной системы ОДУ первого порядка: $\mathbf{y}' = A(x)\mathbf{y}$, где $A(x)$ — матрица размера $n \times n$ с непрерывными элементами.
\textbf{Формула Лиувилля-Остроградского:}
Если $W(x)$ — определитель Вронского фундаментальной системы решений, то:
$$W(x) = W(x_0)\exp\left(\int_{x_0}^x \text{tr}A(t)dt\right)$$
где $\text{tr}A(x) = \sum_{i=1}^n a_{ii}(x)$.
} \hfill \block{
\subsubsection*{10. Формула Лиувилля-Остроградского для ОДУ n-го порядка}
Для линейного однородного дифференциального уравнения n-го порядка:
$$y^{(n)} + p_1(x)y^{(n-1)} + p_2(x)y^{(n-2)} + \ldots + p_{n-1}(x)y' + p_n(x)y = 0$$
\textbf{Формула Лиувилля-Остроградского:}
Если $y_1, y_2, \ldots, y_n$ — фундаментальная система решений и $W(x)$ — их определитель Вронского, то:
$$W(x) = W(x_0)\exp\left(-\int_{x_0}^x p_1(t)dt\right)$$
\textbf{Частные случаи:} Для уравнения второго порядка $y'' + p(x)y' + q(x)y = 0$: $W(x) = W(x_0)\exp\left(-\int_{x_0}^x p(t)dt\right)$.
Для уравнения $y'' + q(x)y = 0$: $W(x) = \text{const}$
\subsubsection*{*Определитель Вронского:}
$$ W(y_1, y_2, y_3) = \begin{vmatrix}
y_1 & y_2 & y_3 \\
y_1' & y_2' & y_3' \\
y_1'' & y_2'' & y_3''
\end{vmatrix}$$
}

\end{document}